\section{Từ trọng số đến bản đồ nhiệt}
\label{sec:ch05-ban-do-nhiet}

\tn{Bản đồ đặc trưng}{feature map} đưa attribution vào không gian có cấu
trúc, nhưng vẫn cần một cách gộp chúng lại thành một hình. Grad-CAM lấy
gradient của điểm số $y^c$ theo từng phần tử
$A^k_{ij}$, rồi trung bình gradient đó trên các vị trí để được trọng số
$\alpha_k^c$:

\begin{equation}
\label{eq:ch05-trong-so-grad-cam}
\alpha_k^c=\frac{1}{Z}\sum_i\sum_j
\frac{\partial y^c}{\partial A^k_{ij}}.
\end{equation}

Ở đây $Z$ là số vị trí trong một bản đồ đặc trưng. Trọng số trong
phương trình~\ref{eq:ch05-trong-so-grad-cam} là
\tn{xấp xỉ tuyến tính}{linear approximation} cục bộ của ảnh hưởng mà bản đồ
đặc trưng $k$ có lên điểm số lớp $c$. Sau đó phương pháp cộng có trọng số các
bản đồ đặc trưng ở lượt truyền xuôi và giữ lại phần dương:

\begin{equation}
\label{eq:ch05-ban-do-grad-cam}
L^c_{\mathrm{Grad\mbox{-}CAM}}=\operatorname{ReLU}
\left(\sum_k \alpha_k^c A^k\right).
\end{equation}

\index{bản đồ nhiệt}%
\tn{Bản đồ nhiệt}{heatmap} $L^c_{\mathrm{Grad\mbox{-}CAM}}$ có cùng độ phân
giải với bản đồ đặc trưng, nên nó được nội suy khi đặt lên ảnh gốc. ReLU trong
phương trình~\ref{eq:ch05-ban-do-grad-cam} không phải là bước trang trí: nó bỏ
các vùng làm giảm điểm số lớp $c$, bởi bản đồ nhiệt mà Grad-CAM dựng lên chỉ
nhằm định vị bằng chứng ủng hộ lớp đang xét. Nếu câu hỏi là vùng phản đối lớp
$c$, cần đổi mục tiêu tính toán thay vì đọc ngược cùng một bản đồ nhiệt.

Hình~\ref{fig:ch05-hai-phep-tich} đặt hai phép gộp của chương này cạnh nhau.
Integrated Gradients lấy tích phân dọc theo một đường trong không gian đầu
vào để chia chênh lệch giữa hai điểm. Grad-CAM trung bình gradient theo không
gian rồi dùng kết quả để gộp các bản đồ đặc trưng. Cả hai đều dùng gradient,
nhưng đối tượng được cộng và câu hỏi được trả lời không giống nhau.

\begin{figure}[htbp]
  \centering
  \begin{tikzpicture}[
  box/.style={draw=black!65, rounded corners=2pt, align=center,
    minimum width=3.1cm, minimum height=0.85cm, fill=black!5},
  data/.style={draw=black!65, rounded corners=2pt, align=center,
    minimum width=3.1cm, minimum height=0.85cm, fill=black!15},
  arrow/.style={-{Stealth[length=2mm]}, draw=black!65, thick},
  nhan/.style={font=\footnotesize, inner sep=1.5pt}
]
  \node[data] (base) {$x_0$: baseline};
  \node[box, right=10mm of base] (path) {Các điểm trên\\đoạn thẳng};
  \node[data, right=10mm of path] (input) {$x$: đầu vào};
  \node[box, below=9mm of path] (ig) {Lấy tích phân gradient\\rồi phân bổ $F(x)-F(x_0)$};
  \draw[arrow] (base) -- (path);
  \draw[arrow] (path) -- (input);
  \draw[arrow] (path) -- node[nhan, right] {phép lấy giới hạn} (ig);

  \node[data, below=21mm of base] (maps) {Bản đồ đặc trưng $A^k$};
  \node[box, right=10mm of maps] (weights) {Gradient của $y^c$\\gộp theo không gian};
  \node[data, right=10mm of weights] (heat) {Bản đồ nhiệt\\cho lớp $c$};
  \draw[arrow] (maps) -- (weights);
  \draw[arrow] (weights) -- (heat);
\end{tikzpicture}

  \caption{Hai phép gộp gradient khác nhau ở chỗ chúng cộng cái gì: Integrated
  Gradients cộng dọc một đường trong không gian đầu vào, còn Grad-CAM cộng
  ngang qua các vị trí của một bản đồ đặc trưng. Mũi tên nét đứt là bước lấy
  tích phân, phần duy nhất không phải một phép cộng hữu hạn. Sơ đồ này diễn
  đọc lại cơ chế của hai bài báo gốc~\autocite{p11ig,p12gradcam}.}
  \label{fig:ch05-hai-phep-tich}
\end{figure}

